\section{Actual holomorphic evaluation of absolutely convergent coefficient series}

Write $P_\rho=\{z:|z_i|<\rho_i\}$ and
$\overline P_\rho=\{z:|z_i|\leq\rho_i\}$ for the original coordinate
polydiscs, and $|F|_\rho=\sum_\alpha|F_\alpha|\rho^\alpha$.
All evaluations below are literal sums of the original complex monomials.

\begin{theorem}[The actual monomial sums converge uniformly and retain the mass bound]
\label{hir:actual-coefficient-polydisc-uniform-evaluation}
\uses{hir:actual-monomial-owner-polydisc-convergence}
\lean{Hironaka.HironakaActualPolydiscHolomorphicEvaluation.closedPolydisc,Hironaka.HironakaActualPolydiscHolomorphicEvaluation.openPolydisc,Hironaka.HironakaActualPolydiscHolomorphicEvaluation.actualSeriesEvaluation,Hironaka.HironakaActualPolydiscHolomorphicEvaluation.evaluation_uniform_and_bound}
If the actual coefficients of $F$ have finite absolute mass at $\rho$,
the finite monomial sums converge uniformly to the actual evaluation on
$\overline P_\rho$, and $|F(z)|\leq|F|_\rho$ there.
\leanok
\end{theorem}
\begin{proof}
Each actual monomial has norm at most $|F_\alpha|\rho^\alpha$
on the entire closed polydisc, by multiplication of the coordinate
inequalities. The given finite mass is precisely summability of these
nonnegative majorants. The uniform Weierstrass test gives the literal
finite-subset sum limit. The norm of the actual sum is at most the sum
of its monomial norms, which is at most the stated mass by the already
proved absolute evaluation estimate.
\end{proof}

\begin{theorem}[Finite actual coefficient mass produces a holomorphic function]
\label{hir:actual-coefficient-polydisc-holomorphy}
\uses{hir:actual-coefficient-polydisc-uniform-evaluation}
\lean{Hironaka.HironakaActualPolydiscHolomorphicEvaluation.evaluation_differentiableOn,Hironaka.HironakaActualPolydiscHolomorphicEvaluation.evaluation_contDiffOn,Hironaka.HironakaActualPolydiscHolomorphicEvaluation.evaluation_zero_neighbourhood}
The actual evaluation is complex differentiable and complex smooth on
$P_\rho$. Positive radii make this a genuine neighbourhood of zero;
no differentiability or convergence of an output function is assumed.
\leanok
\end{theorem}
\begin{proof}
The open polydisc is a finite intersection of inverse images of open
real intervals under continuous coordinate norms. Every actual
monomial is a product of powers of complex coordinate linear maps,
hence is everywhere complex differentiable. To use the genuine
several-variable Weierstrass theorem, encode the countable exponent
type injectively in the natural numbers. Preimages of the finite
initial intervals exhaust every finite exponent set, so their finite
monomial sums retain the proved uniform limit. At any point of the
open polydisc choose a ball contained in it. The finite sums are
complex differentiable on that ball and converge uniformly there.
The several-variable Weierstrass theorem proves differentiability at
its centre. Finally the proved finite-dimensional complex regularity
theorem gives complex smoothness throughout the open polydisc.
\end{proof}

\begin{theorem}[The actual owner quotients and remainder evaluate holomorphically]
\label{hir:actual-monomial-owner-holomorphic-evaluation}
\uses{hir:actual-monomial-owner-absolute-bounds,hir:actual-coefficient-polydisc-holomorphy,hir:actual-coefficient-polydisc-uniform-evaluation}
\lean{Hironaka.HironakaActualPolydiscHolomorphicEvaluation.monomialOwner_evaluations_holomorphic}
For positive radii and finite input mass, every literal actual
monomial-owner quotient and the literal actual remainder give complex
differentiable functions on the same open polydisc. On the entire
closed polydisc their evaluated norms satisfy the actual quotient
bound $\rho^a|Q_aF(z)|\leq|F|_\rho$ and remainder bound
$|RF(z)|\leq|F|_\rho$.
\leanok
\end{theorem}
\begin{proof}
The actual owner-partition estimates internally prove finite mass of
every quotient and remainder, with the stated scaled mass bounds.
Apply the just-proved holomorphy theorem to each unchanged actual
series. The closed-polydisc evaluation bound followed by its actual
mass bound proves the asserted norm inequalities; positive radii
ensure multiplication by $\rho^a$ preserves the inequality.
\end{proof}

This produces genuine holomorphic functions from actual coefficient
summability. It does not assert that an arbitrary holomorphic germ's
formal Taylor series has such a radius, nor prove convergence or
termination of perturbed analytic vertex division or HS tracking.
